% ===== Section 4: "Is this meant to be funny?", humour (~15 min) =====
\topicopener{laughing}{Հումոր}{Սա պիտի որ ծիծաղելի՞ լիներ}{Հումորի տարօրինակ ու անկանխատեսելի աշխատանքը}{happy}

% Dissecting a frog
\begin{frame}{Կատակն ու գորտը}
  \begin{columns}[c]
    \begin{column}{0.54\textwidth}
      \begin{bookquote}
        «Կատակ բացատրելը նման է գորտ հերձելուն. ավելի լավ ես հասկանում, բայց գորտը սատկում է»։\par\smallskip
        {\small Է. Բ. Ուայթ}
      \end{bookquote}
      {\scriptsize Գրքում հեղինակը սխալմամբ «Է. Լ. Ուայթ» է։ Իրականում Է. Բ. Ուայթն ու Քեթրին Ուայթն են, 1941}
    \end{column}
    \begin{column}{0.4\textwidth}
      \onslide<2->{%
      \begin{card}[colback=peach!35!cardwhite]
        \centering\bigemoji[1cm]{microscope}\ \bigemoji[1cm]{laughing}\par\smallskip
        \raggedright Իսկ գիտությունը հենց \textbf{վերլուծելն ու բացատրելն} է։ Գուցե դրա համար են գիտությունն ու հումորը համարվում անհամատեղելի։
      \end{card}}
    \end{column}
  \end{columns}
\end{frame}

% Variables
\begin{frame}{Ինչպես են փորձեր անում}
  \begin{columns}[T]
    \begin{column}{0.47\textwidth}
      \begin{card}[colback=mint!35!cardwhite]
        \emoji[1.3em]{testtube}\ {\blackfont Անկախ փոփոխական} {\scriptsize (independent variable)}\par\smallskip
        Այն, ինչ \textbf{փորձարարն է տալիս}. IQ թեստ, հիշելու բառացանկ
      \end{card}
    \end{column}
    \begin{column}{0.47\textwidth}
      \begin{card}[colback=lavender!35!cardwhite]
        \emoji[1.3em]{chart}\ {\blackfont Կախյալ փոփոխական} {\scriptsize (dependent variable)}\par\smallskip
        Այն, ինչ \textbf{չափում են}. միավորը, հիշած բառերի թիվը, ուղեղի «վառվող» մասերը
      \end{card}
    \end{column}
  \end{columns}
  \vspace{1mm}
  \begin{itemize}
    \item<2-> Անկախ փոփոխականը պետք է \textbf{հուսալիորեն} առաջացնի ցանկալի ռեակցիան
    \item<3-> Ուրեմն հումորի համար պետք է մի բան, որ \textbf{բոլորը} ծիծաղելի կգտնեն
    \item<4-> Ով դա գտնի, երկար գիտնական չի մնա. հեռուստատեսությունը նրան միլիոններ կտա։ Բոլորի սիրած կոմիկ դեռ չի եղել
  \end{itemize}
\end{frame}

% Three more obstacles
\begin{frame}{Եվ ավելի վատ}
  \begin{columns}[T]
    \begin{column}{0.31\textwidth}
      \begin{card}[colback=butter!45!cardwhite]
        \centering\bigemoji[0.9cm]{zany}\par\smallskip
        {\blackfont Անակնկալ}\par\smallskip
        \raggedright Կատակը առաջին անգամ է ծիծաղելի։ Փորձը կրկնելու համար նոր կատակ է պետք
      \end{card}
    \end{column}
    \begin{column}{0.31\textwidth}
      \onslide<2->{%
      \begin{card}[colback=mint!35!cardwhite]
        \centering\bigemoji[0.9cm]{hospital}\par\smallskip
        {\blackfont Լաբորատորիա}\par\smallskip
        \raggedright Ստերիլ, կանոնակարգված միջավայր։ Լավ է գիտության, վատ՝ զվարճանքի համար
      \end{card}}
    \end{column}
    \begin{column}{0.31\textwidth}
      \onslide<3->{%
      \begin{card}[colback=lavender!35!cardwhite]
        \centering\bigemoji[0.9cm]{magnifier}\par\smallskip
        {\blackfont ՄՌՏ (MRI)}\par\smallskip
        \raggedright Նեղ, սառը խողովակ, իսկ հսկա մագնիսը տարօրինակ ձայներ է հանում։ Կատակների տրամադրություն չկա
      \end{card}}
    \end{column}
  \end{columns}
  \vspace{1mm}
  \onslide<4->{\centering\small Ու կրկնելը գիտության մեջ պարտադիր է։ Ուղեղն այնքան անկանխատեսելի է, որ խանգարում է նույնիսկ իրեն ուսումնասիրելուն։\par}
\end{frame}

% Odd strategies: unicycle and the chimney joke
\begin{frame}{Տարօրինակ մեթոդներ}
  \begin{columns}[T]
    \begin{column}{0.47\textwidth}
      \begin{card}[colback=butter!45!cardwhite]
        \emoji[1.3em]{clown}\ {\blackfont Պրոֆեսոր Սեմ Շուստեր}\par\smallskip
        Միանիվ հեծանիվով (unicycle) քշում էր Նյուքասլի փողոցներով ու գրանցում մարդկանց արձագանքները
      \end{card}
      \onslide<2->{%
      \begin{bookquote}
        Բոլորին ծիծաղեցնող բաների առաջին տասնյակում միանիվ հեծանիվը հազիվ թե լինի։
      \end{bookquote}}
    \end{column}
    \begin{column}{0.47\textwidth}
      \onslide<3->{%
      \begin{card}[colback=lavender!35!cardwhite]
        \emoji[1.3em]{speech}\ {\blackfont Պրոֆեսոր Նենսի Բել}\par\smallskip
        Սովորական զրույցների մեջ դիտմամբ \textbf{վատ կատակ} էր խցկում ու գրանցում արձագանքները\par\smallskip
        Կատակն արդեն գիտես։ Այն մեր ծխնելույզներն են \emoji{hourglass}
      \end{card}}
    \end{column}
  \end{columns}
\end{frame}

% The chimney punchline
\begin{frame}{Ի՞նչ ասաց մեծ ծխնելույզը փոքրին}
  \begin{columns}[c]
    \begin{column}{0.56\textwidth}
      \begin{card}[colback=butter!45!cream]
        \centering
        {\blackfont\large Ի՞նչ ասաց մեծ ծխնելույզը\\ փոքր ծխնելույզին}\par\medskip
        \onslide<2->{{\blackfont\LARGE Ոչինչ։}\par\smallskip
        {\blackfont\large Ծխնելույզները չեն խոսում։}}
      \end{card}
      \vspace{2mm}
      \onslide<3->{%
      \begin{chatbox}
        Գնահատիր կատակը \textbf{1-ից 10} (ազնվորեն)
      \end{chatbox}}
    \end{column}
    \begin{column}{0.38\textwidth}
      \onslide<4->{%
      \begin{card}[colback=brainpink!25!cardwhite]
        Արձագանքները տատանվում էին \textbf{անհարմարից} մինչև բացահայտ \textbf{թշնամական}։ Կատակը ոչ ոքի դուր չեկավ, ուստի դա հումորի հետազոտություն համարե՞լ, թե՞ ոչ, վիճելի է։
      \end{card}}
      \onslide<4->{\centering\mascot[1.8cm]{angry}\par}
    \end{column}
  \end{columns}
\end{frame}

% Batman
\begin{frame}{Նենսի Բելի զրուցակիցները}
  \centering
  \begin{meme}[6.2cm]{batman_slapping_robin.jpg}
    \memecap[text width=2.4cm, font=\blackfont\footnotesize]{0.26}{0.87}{Ի՞նչ ասաց մեծ ծխնելույզը փոքրին...}
    \memecap[text width=2.3cm, font=\blackfont\footnotesize]{0.76}{0.89}{Բավակա՛ն է}
  \end{meme}
\end{frame}

% Relief theory
\begin{frame}{Ինչու ենք ծիծաղում. թեթևացում}
  \begin{columns}[c]
    \begin{column}{0.54\textwidth}
      \begin{itemize}
        \item Նիցշե. ծիծաղը մենակության ու մահկանացու լինելու պատասխանն է։ Նրա գրածներից դատելով՝ նա ծիծաղին այդքան էլ ծանոթ չէր
        \item<2-> Ֆրեյդ. ծիծաղը լարվածության \textbf{լիցքաթափում} է։ Այստեղից՝ \textbf{թեթևացման տեսությունը} (relief theory)
        \item<3-> Ուղեղը վտանգ է զգում, այն անվնաս լուծվում է, ու ծիծաղը թուլացնում է լարվածությունը
      \end{itemize}
    \end{column}
    \begin{column}{0.4\textwidth}
      \onslide<4->{%
      \begin{card}[colback=mint!35!cardwhite]
        \emoji[1.4em]{banana}\ Սայթաքեց բանանի կեղևին ու շշմած նստած է. \textbf{ծիծաղելի է}
      \end{card}
      \begin{card}[colback=brainpink!25!cardwhite]
        \emoji[1.4em]{skull}\ Սայթաքեց ու գանգը ճաքեց. \textbf{ծիծաղելի չէ}
      \end{card}}
      \onslide<5->{%
      \begin{card}[colback=lavender!35!cardwhite]
        Հեյվորթ (1928). ծիծաղը նշան է, որ վտանգն \textbf{անցել է}
      \end{card}}
    \end{column}
  \end{columns}
\end{frame}

% Superiority theory
\begin{frame}{Ինչու ենք ծիծաղում. գերազանցություն}
  \begin{columns}[c]
    \begin{column}{0.54\textwidth}
      \begin{itemize}
        \item Դեռ \textbf{Պլատոնն} էր կարծում, որ ծիծաղը գերազանցության արտահայտություն է
        \item<2-> Ինչ-որ մեկն ընկնում է կամ հիմարություն ասում, նրա «կարգավիճակն» իջնում է մերի համեմատ, ու դա մեզ հաճելի է
        \item<3-> Լավ բացատրում է \textbf{չարախնդությունը} (Schadenfreude)
      \end{itemize}
    \end{column}
    \begin{column}{0.4\textwidth}
      \onslide<4->{%
      \begin{card}[colback=lavender!35!cardwhite]
        \centering\bigemoji[1cm]{microphone}\ \bigemoji[1cm]{people}\par\smallskip
        \raggedright Բայց երբ հայտնի կոմիկը մարզադաշտում հազարավոր մարդկանց է ծիծաղեցնում, դժվար թե ամբողջ դահլիճը մտածի՝ «Ինչ հիմար է։ Ես նրանից լավն եմ»։
      \end{card}}
    \end{column}
  \end{columns}
\end{frame}

% Incongruity
\begin{frame}{Ինչու ենք ծիծաղում. խախտված սպասումներ}
  \relax{\small Տեսությունների մեծ մասը շեշտում է \textbf{անհամապատասխանությունը} (incongruity)}\par\medskip
  \centering
  \begin{tikzpicture}[font=\small]
    \node[strip=2.9cm] (s) at (0,0) {\emoji[1.3em]{forkknife}\\ \textbf{Սխեմա} (schema)\\ {\footnotesize ռեստորանում, լողափին, հարցազրույցում ամեն ինչ ընթանում է սովորական ձևով}};
    \node[strip=2.9cm, fill=peach!40!cardwhite, visible on=<2->] (v) at (3.6,0) {\emoji[1.3em]{question}\\ \textbf{Խախտում}\\ {\footnotesize տեղի է ունենում այն, ինչ չպիտի լիներ։ Ուղեղը անորոշություն չի սիրում}};
    \node[strip=2.9cm, fill=mint!40!cardwhite, visible on=<3->] (r) at (7.2,0) {\emoji[1.3em]{bulb}\\ \textbf{Լուծում}\\ {\footnotesize պուանտը (punchline) լուծում է հանելուկը}};
    \node[strip=2.4cm, fill=butter!50!cardwhite, visible on=<4->] (l) at (10.4,0) {\emoji[1.3em]{laughing}\\ \textbf{Ծիծաղ}\\ {\footnotesize ազդանշան՝ «հասկացա»}};
    \draw[arr, visible on=<2->] (s) -- (v);
    \draw[arr, visible on=<3->] (v) -- (r);
    \draw[arr, visible on=<4->] (r) -- (l);
  \end{tikzpicture}\par\medskip
  \raggedright
  \onslide<5->{Դրա համար էլ կրկնված կատակն այդքան ծիծաղելի չէ. ուղեղն արդեն գիտի, որ այն անվնաս է, ու անհամապատասխանություն չկա։\par}
\end{frame}

% Two jokes
\begin{frame}{Երկու կատակ բժշկի ու ձիու մասին}
  \begin{columns}[T]
    \begin{column}{0.47\textwidth}
      \begin{card}[colback=mint!35!cardwhite]
        \emoji[1.3em]{stethoscope}\ «Բժի՛շկ, ինձ թվում է, թե \textbf{վարագույր} եմ»։\par\smallskip
        \onslide<2->{«Դե \textbf{հավաքի՛ր} քեզ»։}
      \end{card}
      \onslide<3->{{\small Ոչ ոք երբեք բժշկի չի գնացել, որովհետև իրեն վարագույր է զգում։ Աբսուրդ իրավիճակ, հետո բառախաղը լուծում է այն։}\par}
    \end{column}
    \begin{column}{0.47\textwidth}
      \onslide<4->{%
      \begin{card}[colback=lavender!35!cardwhite]
        \emoji[1.3em]{beer}\ Ձին մտնում է բար։ Բարմենը հարցնում է. «Why the \textbf{long face}?»
      \end{card}}
      \onslide<5->{{\small Անգլերենում long face նշանակում է «տխուր դեմք», իսկ ձին իսկապես երկար դեմք ունի։ Հայերեն թարգմանված՝ կատակը մեռնում է, ինչպես Ուայթի գորտը։ \textbf{Հումորը հաճախ լեզվի մեջ է}։}\par}
    \end{column}
  \end{columns}
\end{frame}

% Humour in the brain
\begin{frame}{Ծիծաղի կենտրոն չկա}
  \begin{columns}[c]
    \begin{column}{0.56\textwidth}
      \centering
      \begin{tikzpicture}[font=\footnotesize]
        \node[pbox=butter, font=\blackfont\small, minimum size=1.4cm, circle] (h) at (0,0) {\emoji[1.4em]{laughing}};
        \node[strip=3.15cm, visible on=<1->] (a) at (-2.4,1.6) {\textbf{Պարգևի ուղի}\\ ծիծաղը պարգև է};
        \node[strip=3.15cm, visible on=<2->] (b) at (2.4,1.6) {\textbf{Հիպոկամպ, նշաձև մարմին}\\ հիշում ենք՝ ինչ պիտի լիներ};
        \node[strip=3.15cm, visible on=<3->] (c) at (-2.4,-1.6) {\textbf{Ճակատային կեղև}\\ սպասում ու տրամաբանություն};
        \node[strip=3.15cm, visible on=<4->] (d) at (2.4,-1.6) {\textbf{Գագաթային բլթի մասեր}\\ լեզու ու բառախաղ};
        \foreach \n/\s in {a/1, b/2, c/3, d/4}{\draw[plum, line width=1pt, visible on=<\s->] (h) -- (\n);}
      \end{tikzpicture}
    \end{column}
    \begin{column}{0.4\textwidth}
      \onslide<5->{%
      \begin{card}[colback=lavender!35!cardwhite]
        Ըստ մեր այսօրվա գիտելիքի՝ առանձին «ծիծաղի կենտրոն» \textbf{չկա}։ Հումորը բազմաթիվ կապերի արդյունք է, որոնք ձևավորվել են զարգացման, նախասիրությունների ու փորձառությունների ընթացքում։\par\smallskip
        Դրա համար էլ յուրաքանչյուրի հումորը յուրահատուկ է։
      \end{card}}
    \end{column}
  \end{columns}
\end{frame}

% Deaf signers
\begin{frame}{Ժեստերի լեզվով կատակներ}
  \begin{columns}[c]
    \begin{column}{0.54\textwidth}
      \begin{itemize}
        \item Խոսակցության մեջ ծիծաղում ենք \textbf{դադարներին}, նախադասությունների վերջում, որ ծիծաղը չխանգարի լսել
        \item<2-> Ժեստերի լեզվով խոսող խուլ մարդիկ կարող էին ծիծաղել ամբողջ պատմության ընթացքում. ծիծաղի ձայնը ոչինչ չէր ծածկի
        \item<3-> Բայց նրանք ծիծաղում են \textbf{նույն դադարներին}
      \end{itemize}
    \end{column}
    \begin{column}{0.4\textwidth}
      \onslide<4->{%
      \begin{card}[colback=mint!35!cardwhite]
        \centering\bigemoji[1.1cm]{wavehand}\par\smallskip
        \raggedright Լեզուն մշակելու ձևն է որոշում, թե \textbf{երբ} ծիծաղենք։ Ծիծաղն այնքան ինքնաբուխ չէ, որքան կարծում ենք։
      \end{card}}
    \end{column}
  \end{columns}
\end{frame}

% Laughter is social; real vs fake
\begin{frame}{Ծիծաղը սոցիալական է}
  \begin{columns}[c]
    \begin{column}{0.36\textwidth}
      \centering
      % Provine (the book's source): you are about 30 times as likely to laugh with others as alone.
% Alone = 1 as the baseline. xtick listed explicitly (see _learnings pgfplots note).
\begin{tikzpicture}
  \begin{axis}[
      width=5.6cm, height=5.0cm,
      ybar, bar width=1.1cm, bar shift=0pt,
      ymin=0, ymax=36,
      symbolic x coords={alone, group},
      xtick={alone, group},
      xticklabels={{Մենակ}, {Ուրիշների հետ}},
      xticklabel style={align=center, font=\small, text width=2.2cm},
      ytick=\empty, axis line style={plum},
      axis x line*=bottom, axis y line=none,
      nodes near coords={\pgfmathprintnumber{\pgfplotspointmeta}$\times$},
      nodes near coords style={font=\blackfont\normalsize, text=plum},
      tick label style={text=plum},
      enlarge x limits=0.45,
    ]
    \addplot[fill=amBlue, draw=none] coordinates {(alone,1)};
    \addplot[fill=amOrange, draw=none] coordinates {(group,30)};
  \end{axis}
\end{tikzpicture}
\par
      {\scriptsize Ծիծաղելու հավանականությունը (Ռոբերտ Փրովայն)}
    \end{column}
    \begin{column}{0.6\textwidth}
      \begin{itemize}
        \item Ընկերների հետ ծիծաղում ենք նույնիսկ առանց կատակների. ընդհանուր հիշողություններ, սովորական պատմություններ։ Ստենդափն էլ հազվադեպ է մեկ հոգու համար
        \item<2-> Սոֆի Սքոթի փորձերում մարդիկ շատ ճշգրիտ էին տարբերում \textbf{իսկական} ծիծաղը \textbf{կեղծից}
        \item<3-> Դրա համար է սիթքոմների \textbf{ձայնագրված ծիծաղն} այդքան նյարդայնացնում. չենք սիրում, երբ մեր ծիծաղը փորձում են կառավարել
      \end{itemize}
    \end{column}
  \end{columns}
\end{frame}

% When a joke fails
\begin{frame}{Երբ կատակը տապալվում է}
  \begin{columns}[c]
    \begin{column}{0.56\textwidth}
      \begin{itemize}
        \item Կատակ պատմողն իրականում ասում է. «Ես գիտեմ քո հումորը ու կարող եմ քեզ ծիծաղեցնել», այսինքն՝ «կարող եմ քեզ \textbf{կառավարել}»
        \item<2-> Իսկ եթե ուրիշների ներկայությամբ է, իր \textbf{գերազանցությունն} ավելի է շեշտում։ Ուրեմն կատակը պետք է արժենա
        \item<3-> Իսկ եթե չի ծիծաղեցնում, դա ենթագիտակցորեն \textbf{դավաճանություն} է ընկալվում, ու մարդիկ բարկանում են։ Հարցրեք ցանկացած սկսնակ կոմիկի
      \end{itemize}
    \end{column}
    \begin{column}{0.38\textwidth}
      \centering
      \mascot[2.4cm]{sad}\par\medskip
      \onslide<4->{%
      \begin{card}[colback=lavender!35!cardwhite]
        Սա հասկանալու համար պետք է իմանալ, թե ինչպես են \textbf{ուրիշ մարդիկ} ազդում ուղեղի վրա։ Դա արժանի է առանձին գլխի։
      \end{card}}
    \end{column}
  \end{columns}
\end{frame}
